\documentclass{article}

\textwidth=6.5in
\textheight=8.5in
\voffset=-54pt
\oddsidemargin=0in
\evensidemargin=0in
\setlength{\parskip}{8pt}
\setlength{\parindent}{0pt}

\usepackage{workbook}

\usepackage[sols]{handout}

\usepackage{exam}

\newcommand{\answerbox}[3][]{
    \fbox{\begin{minipage}[t][#2][c]{#3}
        #1
    \end{minipage}}
    }

\begin{document} 
 
{\large \mbox{} \hfill \bf Name:\underline{\hspace{3in}}}\\[2pt]
{\mbox{} \hfill \it (Please write your name neatly in print.)}\\[12pt]
{\large \mbox{} \hfill \bf HUID:\underline{\hspace{3in}}}
 
\medskip
 
\noindent
{\large\bf
Exam 2\\
Math Mb Spring 2025}
\noindent


{\Large
\begin{center}\textbf{Directions---Please read carefully!} \end{center}}

This exam is subject to the Harvard College Honor Code. No calculators or any other aids are permitted.  Please write as neatly as possible, as answers that can't be read can't be graded and thus can't receive credit. To receive full credit, you must show relevant working
and include units when appropriate. 

You have 120 minutes to complete this exam. 

You may ask for scratch paper if you need additional space to write, but it will not be scanned or graded. Make sure to include all relevant working in this exam booklet in the space provided for each question.

\begin{center}\textbf{Good luck!}\end{center}


\vfill

\centering{
\emph{As a member of Harvard College, I pledge that I will not give or receive assistance on this exam.}}
\vspace{0.3in}

{\large \mbox{} \hfill \bf Signature:\underline{\hspace{3in}}}
\thispagestyle{empty}

\newpage





\begin{enumerate}


    \item %(SC)
    Evaluate each limit. If the limit does not exist, write DNE, $\infty$, or $-\infty$ as appropriate. 
    
     To receive full credit, you must use correct notation and show all working and/or reasoning. If you use L'Hopital's Rule, you must check that the rule applies and write $\stackrel{\text{L'H}}{=}$ to indicate where you use it.
    \begin{enumerate}
        \item 
        \begin{problem}
        $\displaystyle\lim_{x\to 0}\dfrac{x^2}{\cos x-1}$
        \end{problem}
        \pspace{\vfill}

        \begin{solution}
            We determine the form of the limit by taking the limits of the numerator and denominator separately:
            \[ \lim_{x\to 0} x^2 = 0 \text{ and } \displaystyle \lim_{x\to 0} \cos x -1 = 0
            \]\\
            
            This confirms we have a ``$\dfrac{0}{0}$" form limit, and we can apply L'Hopital's Rule.
            
            In the below computation of the problem, our first application of L'Hopital's Rule gives us another ``$\dfrac{0}{0}$" indeterminate form, so we apply the rule again. Our final limit is not an indeterminate form, and we can evaluate it.
            
            $$\lim_{x\to 0}\dfrac{x^2}{\cos x-1} \stackrel{\text{L'H}}{=} \lim_{x\to 0} \dfrac{2x}{-\sin x}\stackrel{\text{L'H}}{=}\lim_{x\to 0} \frac{2}{-\cos x} = \fbox{-2}$$
        \end{solution}
        \item 
        \begin{problem}
        $\displaystyle\lim_{x\to \infty}\dfrac{x}{\arctan(x^2)}$
        \end{problem}
        \pspace{\vfill}
        
        \begin{solution}
            We determine the form of the limit by taking the limits of the numerator and denominator separately:
            \[
            \lim_{x\to \infty} x = \infty \text{ and } 
            \displaystyle \lim_{x\to \infty} \arctan(x^2) = \dfrac{\pi}{2}
            \]
            This is not an indeterminate form. A numerator growing without bound over a denominator approaching a finite value means that our limit grows without bound.\\

            \[\lim_{x \to \infty} \frac{x}{\arctan(x^2)} = \fbox{$\infty$}\]
            
            
        \end{solution}
        \item 
        \begin{problem}
        $\displaystyle\lim_{x\to -\infty}\dfrac{e^x+10x+4}{x-1}$
        \end{problem}
        \pspace{\vfill}

        \begin{solution}
            We determine the form of the limit by taking the limits of the numerator and the denominator separately: \\
            $\displaystyle \lim_{x\to-\infty} e^x + 10x + 4 =  - \infty$ because $e^x \to 0$ and $10x\to -\infty$ and 
            $\displaystyle \lim_{x\to-\infty} x-1 = -\infty$\\
            This is an indeterminate form $``\frac{\infty}{\infty}$" \\
            
            There are multiple options here. Option A, we could apply L'Hopital's Rule once:
            $$\lim_{x\to -\infty}\dfrac{e^x+10x+4}{x-1} \stackrel{\text{L'H}}{=} \lim_{x\to-\infty} \frac{e^x+10}{1} = \fbox{10}$$
            Option B, we can use relative growth rates to identify the fastest growing terms. As $x\to -\infty$, $10x>>e^x$ is the fastest growing term in the numerator (since $e^x$ goes to 0). $x$ is the fastest growing term in the denominator. Having established that relative growth behavior, we can compute the limit:
            $$\lim_{x\to -\infty}\dfrac{e^x+10x+4}{x-1}=\lim_{x\to-\infty}\frac{10x}{x} = \fbox{10}$$
        \end{solution}

        % add a few more?
        
    \end{enumerate}

\pspace{\newpage}

\probonly{
\begin{framed}
        \emph{Continued from the previous page. The problem statement is provided again below}
    
Evaluate each limit. If the limit does not exist, write DNE, $\infty$, or $-\infty$ as appropriate. 
    
     To receive full credit, you must use correct notation and show all working and/or reasoning. If you use L'Hopital's Rule, you must check that the rule applies and write $\stackrel{\text{L'H}}{=}$ to indicate where you use it.

    \end{framed}
    }

\begin{enumerate}[resume]

        
        \item 
        \begin{problem}
        $\displaystyle \lim_{x\to\infty}\sqrt{x} \sin(\tfrac{1}{x})$
        \end{problem}
        \pspace{\vfill}

        \begin{solution}
            We determine the form of the limit by taking the limits of the terms in the product individually:\\
            $\displaystyle \lim_{x\to \infty} \sqrt{x} = \infty$ and $\displaystyle \lim_{x\to\infty}\sin\bigg(\frac{1}{x}\bigg) = 0$.\\
            
            This is a $``0\cdot \infty"$ indeterminate form. We can express this as a quotient to use L'Hopital's Rule. (Note: In the work below, the first equality represents a rewriting of the functions so that taking a derivative is clearer. The second equality is the result of L'Hopital's Rule, and the 3rd equality is a result of algebraically simplifying. The final limit is no longer indeterminate. The numerator approaches 2 and the denominator grows without bound, so it can be be computed directly.
            $$  \lim_{x\to\infty}\sqrt{x} \sin(\frac{1}{x}) = \lim_{x\to\infty}\frac{\sin(x^{-1})}{x^{-1/2}}\stackrel{\text{L'H}}{=}\lim_{x\to\infty}\frac{\cos(x^{-1})\cdot (-x^{-2})}{(\frac{-1}{2})x^{-3/2}}=\lim_{x\to\infty}\frac{2\cos(\frac{1}{x})}{x^{1/2}}=\fbox{0}$$
        \end{solution}
\end{enumerate}

\item 
\begin{problem}
    Find the derivative of each function.
\end{problem}

\begin{enumerate}
    \item 
    \begin{problem}
    $A(x)=\displaystyle \int_0^x \arcsin(t)\,dt$
    \end{problem}
    \pspace{\stretch{0.4}}

\begin{solution}
    By the Fundamental Theorem of Calculus, we know that an area function defined in this way will be changing at a rate of \fbox{$\arcsin(x)$}.
\end{solution}
    \item 
    \begin{problem}
    $B(x)=\displaystyle \int_0^x e^t\,dt -\int_2^x e^t\,dt$
    \end{problem}
    \pspace{\stretch{0.4}} 
    \begin{solution}
        Through properties of the definite integral, we can simplify this into: 
        $\displaystyle \int_0^x e^t\,dt -\int_2^x e^t\,dt = \int_0^x e^t\,dt +\int_x^2 e^t\,dt = \int_0^2 e^t \, dt$
        This is evaluates to a constant. The derivative of a constant is \fbox{0}
    \end{solution}
    
\end{enumerate}

\pspace{\newpage}

\item % SC
\begin{problem}
 The MBTA is testing the acceleration of their new train cars. Let $v(t)$ be the car's speed in feet per second $t$ seconds after the car starts moving. The table below shows the readings on the car's speedometer every 2 seconds as the car's speed continuously increases.

    
    \begin{tabular}{c|c|c|c|c|c}
        $t$ (seconds) &  0 & 2 & 4 & 6 & 8 \\
        \hline
        $v(t)$ (feet per second) & 0 & 25 & 45 & 60 & 70
    \end{tabular}

 \end{problem}
    \begin{enumerate}
        \item 
        \begin{problem}
        Using correct units, explain the meaning of $\displaystyle \int_0^8 v(t)\,dt$ in the context of this problem.
        \end{problem}
        \pspace{\vfill}

        \begin{solution}
            Since $v(t)$ is speed, and speed is the rate function for distance traveled, this definite integral will represent the total distance traveled by the MBTA car (in feet) during the first 8 seconds after the car begins moving.
        \end{solution}
        \item 
        \begin{problem}
        Approximate $\displaystyle \int_0^8 v(t)\,dt$ with a left Riemann sum with 4 equal sub-intervals, $L_4$. Include units.
        \end{problem}
        \pspace{\vfill}

\begin{solution}
When calculating an approximation using 4 subintervals over an 8 second interval, we will want to use 2 second subintervals. Each subinterval will use the speed at the $t$-value on the left side of the subinterval, that is, at $t=0, 2, 4, \text{ and } 6$
    \[
    L_4= 0(2)+25(2)+45(2)+60(2) = 50+90+120=\fbox{$260 \text{ feet}$}
    \]
\end{solution}
        This is an (circle one) \hfill \textbf{underestimate} \hfill \textbf{overestimate} \hfill \textbf{not enough information} \hfill

        \begin{solution} 
        We are told in the problem prompt that the speed is continuously increasing. This means that using the speed at the lefthand side (or start) of each subinterval will be an underestimate for its speed over the rest of the interval, so the approximate distance traveled will also be an \fbox{underestimate}.
        \end{solution} 
        
        \item 
        \begin{problem}
            Approximate $\displaystyle \int_0^8 v(t)\,dt$ with a right Riemann sum with 4 equal sub-intervals, $R_4$. Include units.
        \end{problem}
        \pspace{\vfill}
        \begin{solution}
            Same as before, when calculating an approximation using 4 subintervals over an 8 second interval, we will want to use 2 second subintervals. This time, each subinterval will use the speed at the $t$-value on the right side (or end) of the subinterval, that is, at $t= 2, 4, 6,\text{ and } 8$
    \[
    R_4= 25(2)+45(2)+60(2) + 70(2) = 50+90+120 + 140=\fbox{$400 \text{ feet}$}
    \]
        \end{solution}

        This is an (circle one) \hfill \textbf{underestimate} \hfill \textbf{overestimate} \hfill \textbf{not enough information} \hfill

\begin{solution}
    We return to the fact that the speed is continuously increasing. This means that using the speed at the righthand side (or start) of each subinterval will be an overestimate for its speed over the rest of the interval, so the approximate distance traveled will also be an \fbox{overestimate}.
\end{solution}
        \item 
\begin{problem}
Suppose we were also able to compute $R_{8}$, $L_{8}$, $R_{80}$, and $L_{80}$. Order these, along with  $\displaystyle\int_0^{8} v(t)\,dt$, from \textbf{least to greatest.}
\end{problem}

    \newcommand{\blank}{\quad\underline{\hspace{2cm}}\quad}

\probonly{
    \[\hss
    \blank< \blank < \blank < \blank < \blank
    \hss\]
}

\begin{solution}
    The lefthand approximations are all underestimates. The righthand approximations are all overestimates. Increasing the number of intervals will improve the approximation and bring it closer to the actual value of the definite integral. This gives us:
    \[
    L_8 < L_{80} < \int_0^8 v(t) dt<R_{80} < R_8
    \]
\end{solution}
    \end{enumerate}


\pspace{\newpage}

    \item %(YH)
    Let $f$ be a function where $\displaystyle \int_1^4f(t)\,dt= 6$. Evaluate each expression below or say that there is not enough information to do so. Show the working or reasoning that leads to your answer.
    \begin{enumerate}
    % cut some of these?
        % \item 
        % \begin{problem}
        % $\displaystyle\int_2^5 4f(t)\,dt$
        % \end{problem}
        % \pspace{\vfill}
        
        % \item 
        % \begin{problem}
        % $\displaystyle \int_2^5 f(t)\,dt+9$
        % \end{problem}
        % \pspace{\vfill}
        
        \item 
        \begin{problem}
        $\displaystyle \int_4^1 \big(2f(t)+10\big)dt$
        \end{problem}
        \pspace{\vfill}

        \begin{solution}
           \begin{align*}
           \int_4^1 \big(2f(t)+10\big)\,dt & =\int_4^12 f(t)\,dt + \int_4^1 10\,dt \\
           &= -2\int_1^4f(t)\,dt - \int_1^4 10\,dt \\
           &= -2(6)-30 \\ &= \boxed{-42}
           \end{align*}
           Note that the first equality is justified that the integral of the sum of two functions is equal to the sum of the function's individual integrals. The second equality
           
        \end{solution}
        \item 
        \begin{problem}
        $\displaystyle\int_1^4 f(t+5)\,dt$
        \end{problem}
        \pspace{\vfill}

\begin{solution}
    The function $f(t+5)$ is a transformation of $f(t)$ where the graph has been horizontally shifted 5 units to the left, so considering this graphically will help us. Our given information told us about the signed area under the curve $f(t)$ on the interval $[1,4]$. We know no other information about the function. If we shift the graph to the left by 5 units, the interval containing the known signed area has also shifted to the left, and now exists over the interval $[-4, -1]$ (we shifted both endpoints to the left by 5 as well). The signed area of our shifted function without also shifting the interval is unknown, so there is \fbox{not enough information}. 
\end{solution}
        \item 
        \begin{problem}
        $\displaystyle\int_6^9 f(t+5)\,dt$
        \end{problem}
        \pspace{\vfill}

        \begin{solution}
            There is \fbox{not enough information} to compute this integral. Using the same reasoning as in part (b), when the function shifts to the left by 5 units, the interval where signed area is known also shifts to the left by 5 units. This definite integral has shifted the interval to the right instead of the left. 
        \end{solution}
        % \item 
        % \begin{problem}
        % $\displaystyle\int_0^{-3} f(t+5)\,dt$
        % \end{problem}
        % \pspace{\vfill}

        \end{enumerate}
\pspace{\newpage}

\probonly{
\begin{framed}
        \emph{Continued from the previous page. The problem statement is provided again below}
    
Let $f$ be a function where $\displaystyle \int_1^4f(t)\,dt= 6$. Evaluate each expression below or say that there is not enough information to do so. Show the working or reasoning that leads to your answer.
    \end{framed}
    }
    
\begin{enumerate}[resume]

        \item
        \begin{problem}
            $\displaystyle \int_1^4 \big(f(t)-2\pi t)\,dt$
        \end{problem}
        \pspace{\vfill}

\begin{solution}
    We can first isolate our known integral accordingly: \[\displaystyle \int_1^4 \big(f(t)-2\pi t)\,dt = \int_1^4 f(t) \, dt - 2\pi \int_1^4 t\, dt\]
    
    On the righthand side, the former integral is our given information. The latter integral can be reasoned geometrically. Sketching the graph of the function $t$, the signed area under the curve makes a trapezoid, whose area is $\frac{1}{2}(1+4)(3)=\frac{15}{2}$
   \[\int_1^4 \big(f(t)-2\pi t)\,dt = \int_1^4 f(t) \, dt - 2\pi \int_1^4 t\, dt\ = 6 - 2\pi \cdot \frac{15}{2} = \fbox{$6-15\pi$}\]
\end{solution}
        
        \item 
        \begin{problem}
        $\displaystyle
        \int_1^0 f(t)\,dt + \int_0^4 \big(f(t)+\sqrt{16-t^2}\big)\,dt$
        \end{problem}
        \pspace{\stretch{1.5}}

        \begin{solution}
        We first attempt to isolate definite integrals involving $f(t)$ and try to combine them, in the hope of identifying our given information in the expression.
         \begin{align*}
         \int_1^0 f(t)\,dt + \int_0^4 \big(f(t)+\sqrt{16-t^2}\big)\,dt &= -\int_0^1f(t) \, dt + \int_0^4 f(t) \, dt + \int_0^4 \sqrt{16-t^2}\, dt \\
         &= \int_1^4 f(t) \, dt+ \int_0^4 \sqrt{16-t^2}\, dt
         \end{align*}
         The former integral is our given information. The latter integral can be reasoned geometrically, the function is a semicircle with radius 4 and the interval refers to the signed area on the right half of it. Its value is one quarter the area of a circle with radius 4.
         \[\int_1^4 f(t) \, dt+ \int_0^4 \sqrt{16-t^2}\, dt = 6 + \frac{\pi\cdot (4^2)}{4} = \fbox{$6+4\pi$}\]
        \end{solution}
    \end{enumerate}

\pspace{\newpage}

% \item
% \begin{problem}
%     $f$ is a continuous piecewise-linear function defined on $[-3,6]$. The graph of $y=f(x)$ is shown below.
% \end{problem}

% \vspace{-12pt}
% \begin{center}
%   \begin{tikzpicture}
%     \begin{axis}[xtick={-5,-4, -3, -2, -1, 0, 1, ..., 6,7}, ytick={-3,-2, -1, ..., 5},
%     xlabel=$x$, ylabel=$y$,
%       grid=major, cycle list=black, axis equal=true, ymin=-3, ymax=5,
%       clip=false, width=4in, height=3.05in, axis line style={shorten >=-8pt}, y label style={rotate=-90, yshift=8pt}, x label style={xshift=8pt}]
%       \addplot[closed] coordinates {(-3,-2) (0,4) (4,0) (6,4)};
%       \addplot+[sharp plot] coordinates { (-3, -2) (0,4) (4, 0) (6, 4) };
%     \end{axis}
%   \end{tikzpicture}
%   \end{center}
%   \vspace{-12pt}
% The function $F$ is defined on the domain $[-3, 6]$ as $F(x)=\displaystyle \int_0^x f(t)\,dt$.

% \begin{enumerate}
%     \item
%     \begin{problem}
%         Evaluate $F(-3)$ and $F(6)$.
%     \end{problem}
%     \pspace{\vfill}
    
%     \item 
%     \begin{problem}
%         Find \textbf{all} $x$-values for which $F(x)=0$. Justify your answer briefly.
%     \end{problem}
%     \pspace{\vfill}


% \end{enumerate}
% \pspace{\newpage}

% \probonly{
% \begin{framed}
%         \emph{Continued from the previous page. The problem statement is provided again below}
    
% $f$ is a continuous piecewise linear function defined on $[-3,6]$. The graph of $y=f(x)$ is shown below.

% \vspace{-12pt}
% \begin{center}
%   \begin{tikzpicture}
%     \begin{axis}[xtick={-5,-4, -3, -2, -1, 0, 1, ..., 6,7}, ytick={-3,-2, -1, ..., 5},
%     xlabel=$x$, ylabel=$y$,
%       grid=major, cycle list=black, axis equal=true, ymin=-3, ymax=5,
%       clip=false, width=4in, height=3.05in, axis line style={shorten >=-8pt}, y label style={rotate=-90, yshift=8pt}, x label style={xshift=8pt}]
%       \addplot[closed] coordinates {(-3,-2) (0,4) (4,0) (6,4)};
%       \addplot+[sharp plot] coordinates { (-3, -2) (0,4) (4, 0) (6, 4) };
%     \end{axis}
%   \end{tikzpicture}
%   \end{center}
%   \vspace{-12pt}
% The function $F$ is defined on the domain $[-3, 6]$ as $F(x)=\displaystyle \int_0^x f(t)\,dt$.
%     \end{framed}
%     }
    
% \begin{enumerate}[resume]

%         \item 
%     \begin{problem}
%         Determine the interval(s) on which $F$ is increasing. Justify your answer briefly.
%     \end{problem}
%     \pspace{\vfill}
    
%     \item 
%     \begin{problem}
%         Determine the interval(s) on which $F$ is concave up. Justify your answer briefly.
%     \end{problem}
%     \pspace{\vfill}

% \end{enumerate}
% \pspace{\newpage}

% \probonly{
% \begin{framed}
%         \emph{Continued from the previous page. The problem statement is provided again below}
    
% $f$ is a continuous piecewise linear function defined on $[-3,6]$. The graph of $y=f(x)$ is shown below.

% \vspace{-12pt}
% \begin{center}
%   \begin{tikzpicture}
%     \begin{axis}[xtick={-5,-4, -3, -2, -1, 0, 1, ..., 6,7}, ytick={-3,-2, -1, ..., 5},
%     xlabel=$x$, ylabel=$y$,
%       grid=major, cycle list=black, axis equal=true, ymin=-3, ymax=5,
%       clip=false, width=4in, height=3.05in, axis line style={shorten >=-8pt}, y label style={rotate=-90, yshift=8pt}, x label style={xshift=8pt}]
%       \addplot[closed] coordinates {(-3,-2) (0,4) (4,0) (6,4)};
%       \addplot+[sharp plot] coordinates { (-3, -2) (0,4) (4, 0) (6, 4) };
%     \end{axis}
%   \end{tikzpicture}
%   \end{center}
%   \vspace{-12pt}
% The function $F$ is defined on the domain $[-3, 6]$ as $F(x)=\displaystyle \int_0^x f(t)\,dt$.
%     \end{framed}
%     }
    
% \begin{enumerate}[resume]


% \item 
% \begin{problem}
%     Find the absolute maximum and minimum values of $F(x)$ on $[-3,6]$. Give also the $x$-values where the maximum and minimum are attained. Justify your answer fully.
% \end{problem}
    
% \end{enumerate}

% \pspace{\newpage}

\item 
\begin{problem}
    A squirrel is climbing straight up and down a tall tree in Harvard Yard. At time $t=0$, the squirrel is 5 feet above the ground. For $0\leq t\leq 18$, the squirrel's velocity $v(t)$ (in feet per second) at time $t$ (in seconds) is modeled by the piecewise-linear function defined by the graph shown.   
\vspace{-6pt}
\begin{center}
  \begin{tikzpicture}
    \begin{axis}[xtick={0,2,...,18}, minor x tick num=1, 
    ytick={-2,0, ..., 4}, minor y tick num=1,
    xlabel=$t$ (s), ylabel=$v(t)$ (ft/s),
    cycle list=black, axis equal=true, ymin=-3, ymax=5, xmin=0, xmax=19,
      clip=false, width=4in, height=2in]
      \addplot[closed] coordinates {(0,0) (2,4) (7,4) (10,-2) (14,-2) (16,2) (18,2)};
      \node[above] at (2,4) {$(2,4)$};
      \node[above] at (7,4) {$(7,4)$};
      \node[below] at (10,-2) {$(10,-2)$};
      \node[below] at (14,-2) {$(14,-2)$};
      \node[above] at (16,2) {$(16,2)$};
       \node[below] at (18,2) {$(18,2)$};
      \addplot+[sharp plot] coordinates {(0,0) (2,4) (7,4) (10,-2) (14,-2) (16,2) (18,2)};
    \end{axis}
  \end{tikzpicture}
  \end{center}
  Answer each part based on the model. Show all work and/or reasoning and include units.
 \end{problem}

\begin{enumerate}
    % \item 
    % \begin{problem}
    %     At what times in the interval $0<t<18$, if any, does the squirrel change direction? Explain.
    % \end{problem}
    % \pspace{\vfill}

    \item 
    \begin{problem}
        At what time in the interval $0\leq t \leq 18$ was the squirrel highest above the ground? How high above the ground is the squirrel at this time? 
    \end{problem}
    \pspace{\stretch{2}}
    
\begin{solution}
    The squirrel is highest above the ground at \fbox{$t=9$}. When $v(t)$ is positive, this means the squirrel is moving up the tree and when $v(t)$ is negative, the squirrel is moving down the tree. The signed area under the curve gives the squirrel's net change in height off the ground. The most substantial accumulation in height ends at $t=9$ before the squirrel begins moving downward. Though the squirrel descends and ascends again, the decrease in the squirrel's height on the interval $[9, 15]$ is larger than the final increase $[15, 18]$. The squirrel's height is calculated as follows:
    \[5+\int_0^9 v(t)\, dt = 5+\frac{1}{2}(9+5)(4) = 5+28=\fbox{$33$ feet}\]
\end{solution}
%     \end{enumerate}
    
% \pspace{\newpage}

% \probonly{
% \begin{framed}
%         \emph{Continued from the previous page. The problem statement is provided again below.}
    
% A squirrel is climbing straight up and down a tall tree in Harvard Yard. At time $t=0$, the squirrel is 5 feet above the ground. For $0\leq t\leq 18$, the squirrel's velocity $v(t)$ (in feet per second) at time $t$ (in seconds) is modeled by the piecewise-linear function defined by the graph shown.   
% \vspace{-6pt}
% \begin{center}
%   \begin{tikzpicture}
%     \begin{axis}[xtick={0,2,...,18}, minor x tick num=1, 
%     ytick={-2,0, ..., 4}, minor y tick num=1,
%     xlabel=$t$ (s), ylabel=$v(t)$ (ft/s),
%     cycle list=black, axis equal=true, ymin=-3, ymax=5, xmin=0, xmax=19,
%       clip=false, width=4in, height=2in, y label style={rotate=-90}]
%       \addplot[closed] coordinates {(0,0) (2,4) (7,4) (10,-2) (14,-2) (16,2) (18,2)};
%       \node[above] at (2,4) {$(2,4)$};
%       \node[above] at (7,4) {$(7,4)$};
%       \node[below] at (10,-2) {$(10,-2)$};
%       \node[below] at (14,-2) {$(14,-2)$};
%       \node[above] at (16,2) {$(16,2)$};
%        \node[below] at (18,2) {$(18,2)$};
%       \addplot+[sharp plot] coordinates {(0,0) (2,4) (7,4) (10,-2) (14,-2) (16,2) (18,2)};
%     \end{axis}
%   \end{tikzpicture}
%   \end{center}
%     \end{framed}
%     }
    
% \begin{enumerate}[resume]

\item
\begin{problem}
    Find the total distance the squirrel travels during the interval $0\leq t \leq 18$.
\end{problem}
\pspace{\stretch{2}}

\begin{solution} 
Total distance traveled is found by summing all of the net changes in height without taking sign into account. That is, we want to count feet that the squirrel ascended and descended together in this total. We will calculate this geometrically, adding together the areas between $v(t)$ and the $t$-axis with all terms positive. (Note that work from part a can be reused for the first trapezoid):
\[ \int_0^{18}|v(t)|\,dt = 28+\frac{1}{2}(6+4)(2) + \frac{1}{2}(3+2)(2) = 28+10+5=\fbox{43 feet}\]
\end{solution}



% \item
% \begin{problem}
%     Write a formula involving a definite integral for the squirrel's height above the ground $h(t)$ (in feet) at time $t$ (in seconds).
% \end{problem}
% \pspace{\vfill}
\end{enumerate}

\pspace{\newpage}

\item 
\begin{problem}
    At Bagelsaurus, a famous Cambridge bagel shop, bagels are cooked by first being boiled in a large water bath. Unfortunately, as the water bath was being filled this morning, it started leaking! The bath was being filled at a rate modeled by $f(t)$ liters per minute, and water was leaking out at a rate modeled by $g(t)$ liters per minute. The graphs of the functions $f$ and $g$ are shown below for $0\leq t \leq 8$, where $t$ is measured in minutes. At time $t=0$, there were $10$ liters of water in the bath.

\vspace{-12pt}
    \begin{center}
        \begin{tikzpicture}
            \begin{axis}[xmin=0, xmax=8.5, ymin=0,
            clip=false, grid=major,
            xlabel=$t$ (min), ylabel=rate (L/min), width=3in]
                \addplot[densely dashed, domain=0:8] {20*sin(deg(x^2/35))} node[right] {$g$};
                \addplot[very thick, domain=0:8] {-0.04*x^3+0.4*x^2+1.25*x} node[right] {$f$};
            \end{axis}
        \end{tikzpicture}
    \end{center}
    \vspace{-18pt}
\end{problem}

\begin{enumerate}

    \item 
    \begin{problem}
        Was the amount of water in the bath increasing or decreasing at time $t=2$ minutes? Explain.
    \end{problem}
    \pspace{\vfill}

\begin{solution}
    The amount of water in the bath was increasing at $t=2$ because $f(2)>g(2)$.
\end{solution}
    \item 
    \begin{problem}
        Write an expression involving a definite integral that, when evaluated, would give the amount of water that leaked out of the bath during the 8-minute time interval $0\leq t \leq 8$.
    \end{problem}
    \pspace{\vfill}

    \begin{solution}
        \[\fbox{$\displaystyle \int_0^8 g(t)\, dt$}\]
    \end{solution}

    \item 
    \begin{problem}
        Write an expression involving a definite integral that, when evaluated, would give the total amount of water in the bath at time $t=8$ minutes.
    \end{problem}
    \pspace{\vfill}

    \begin{solution}
        \[\fbox{$\displaystyle 10+ \int_0^8 (f(t) -g(t))\, dt$}\]
    \end{solution}
\end{enumerate}
\pspace{\newpage}

\probonly{
\begin{framed}
        \emph{Continued from the previous page. The problem statement is provided again below}
    
At Bagelsaurus, a famous Cambridge bagel shop, bagels are cooked by first being boiled in a large water bath. Unfortunately, as the water bath was being filled this morning, it started leaking! The bath was being filled at a rate modeled by $f(t)$ liters per minute, and water was leaking out at a rate modeled by $g(t)$ liters per minute. The graphs of the functions $f$ and $g$ are shown below for $0\leq t \leq 8$, where $t$ is measured in minutes. At time $t=0$, there were $10$ liters of water in the bath.

\vspace{-12pt}
    \begin{center}
        \begin{tikzpicture}
            \begin{axis}[xmin=0, xmax=8.5, ymin=0,
            clip=false, grid=major,
            xlabel=$t$ (min), ylabel=rate (L/min), width=3in]
                \addplot[densely dashed, domain=0:8] {20*sin(deg(x^2/35))} node[right] {$g$};
                \addplot[very thick, domain=0:8] {-0.04*x^3+0.4*x^2+1.25*x} node[right] {$f$};
            \end{axis}
        \end{tikzpicture}
    \end{center}
    \end{framed}
    }
    
\begin{enumerate}[resume]
\item 
\begin{problem}
    At what time $t$, $0 \leq t \leq 8$, was the amount of water in the bath the \textbf{highest}? Justify your answer.
\end{problem}
\pspace{\vfill}

\begin{solution}
    The amount of water in the bath was the highest at $t=4$. On the interval $(0,4)$, $f(t)>g(t)$, so the amount of water in the bath is increasing. On the interval (4,8), $f(t)<g(t)$, so the amount of water in the bath is decreasing. This makes $t=4$ the moment when the greatest amount of water will be in the bath.
\end{solution}

\item 
\begin{problem}
    At what time $t$, $0 \leq t \leq 8$, was the amount of water in the bath the \textbf{lowest}? Give an answer, then choose the combination of statements that gives a complete and correct justification.
\end{problem}
\pspace{\stretch{0.25}}

\begin{solution}
    The amount of water in the bath is lowest at $t=8$. The correct justification below is $\fbox{(D)}$
\end{solution}
    \begin{enumerate}[label=\Roman*.]
      \item Because the amount of water increased until $t=4$ then decreased afterwards.
      \item Because the amount of water decreased until $t=4$ then increased afterwards.
      \item Because the net change in the amount of water from $t=0$ to $t=8$ was positive.
      \item Because the net change in the amount of water from $t=0$ to $t=8$ was negative.
    \end{enumerate}

    \begin{enumerate*}[label=(\Alph*), itemjoin=\hfill]
    \item II only
    \item I and III only
    \item II and III only
    \item I and IV only
    \end{enumerate*}
    



\end{enumerate}

\pspace{\newpage}

\item
\begin{problem}
    Shown below is the graph of $f(x)=x\sin x$ on the interval $[0,2\pi]$.

    \begin{center}
        \begin{tikzpicture}
            \begin{axis}[
            xlabel=$x$, ylabel=$y$,
            xtick={0, 1.57, ..., 6.281}, ytick=\empty,
            xticklabels={$0$, $\frac{\pi}{2}$, $\pi$, $\frac{3\pi}{2}$, $2\pi$}, width=3in, clip=false
            ]
                \addplot[domain=0:2*pi] {x*sin(deg(x))};
                \addplot[draw=none, domain=0:6.5] {0};
                \node at (2,2.3) {$y=x\sin x$};
            \end{axis}
        \end{tikzpicture}
    \end{center}   
    \pspace{-12pt}
\end{problem}

\begin{enumerate}
    \item 
    \begin{problem}
        Approximate $\displaystyle \int_0^{2\pi} x \sin x\,dx$ using a right Riemann sum with 4 equal sub-intervals, $R_4$. Show the work that leads to your answer and simplify fully.
    \end{problem}
    \pspace{\vfill}

    \begin{solution}
        The interval $[0, 2\pi]$ split into 4 equal subintervals means the width of each interval will be $\pi/2$. Our approximation will be calculated using the right-hand endpoint of each subinterval:
        \begin{align*}
            R_4 &= \tfrac{\pi}{2}f(\tfrac{\pi}{2}) + \tfrac{\pi}{2}f(\pi)+\tfrac{\pi}{2}f(\tfrac{3\pi}{2}) + \tfrac{\pi}{2}f(2\pi)\\
            &=\tfrac{\pi}{2}(\tfrac{\pi}{2}\sin\tfrac{\pi}{2}) + \tfrac{\pi}{2}(\pi\sin\pi)+
            \tfrac{\pi}{2}(\tfrac{3\pi}{2}\sin\tfrac{3\pi}{2})+
            \tfrac{\pi}{2}(2\pi\sin2\pi)\\
            &=\tfrac{\pi^2}{4}-\tfrac{3\pi^2}{4}\\
            &=\boxed{\frac{-\pi^2}{2}}
        \end{align*}
    \end{solution}

    \item
    \begin{problem}
        Which of the following expressions is equal to $\displaystyle \int_0^{2\pi} x \sin x\,dx$? Circle the letter for all that apply.
    \end{problem}

    
    \begin{enumerate}[label=\protect\maybecircle{\Alph*.}]
      \plainitem $\displaystyle \lim_{n\to\infty} \sum_{k=1}^n \sin(\tfrac{2\pi k}{n})$
      \plainitem $\displaystyle \lim_{n\to\infty} \sum_{k=1}^n\sin(\tfrac{2\pi k}{n})\tfrac{2\pi}{n}$
      \plainitem $\displaystyle \lim_{n\to\infty} \sum_{k=1}^n\tfrac{2\pi k}{n}\sin(\tfrac{2\pi k}{n})$
      \circleitem $\displaystyle \lim_{n\to\infty} \sum_{k=1}^{n}\tfrac{2\pi k}{n}\sin(\tfrac{2\pi k}{n})\tfrac{2\pi}{n}$
      \plainitem None of the above \vphantom{$\displaystyle\sum_{k=0}^{n-1}$}
      \end{enumerate}

\end{enumerate}

\begin{solution}
    By the definition of the definite integral, the integral is equal to $\lim_{n\to\infty} R_n$. Since each computes the appropriate limit, we need to determine which correctly calculates $R_n$. 
    \[R_n = \sum_{k=1}^n f(x_k)\Delta x\]
    In this problem, $\displaystyle \Delta x = \frac{2\pi-0}{n}=\frac{2\pi}{n}$. We also know $\displaystyle x_k = x_0 + k\Delta x = 0+k\Delta x = 0 + k\frac{2\pi}{n} = \frac{2\pi k}{n}$
    Substituting these in for $\Delta x$ and $x_k$, we have:
     \[R_n = \sum_{k=1}^n f(x_k)\Delta x= \sum_{k=1}^n x_k\sin(x_k)\Delta x = \sum_{k=1}^n \frac{2\pi k}{n} \sin\bigg(\frac{2\pi k}{n}\bigg)\cdot \frac{2\pi}{n}\]
     This matches the sigma notation in D.
\end{solution}
\end{enumerate}

\probonly{
\newpage
\textit{This page has been intentionally left blank for scratch work. No work on this page will be graded.}
}

\probonly{
\newpage
\textit{This page has been intentionally left blank for scratch work. No work on this page will be graded.}
}
\end{document} 